\section{The augmentation code on 253 heptads}

Let $\mathcal D$ be the specified $S(4,7,23)$, with point set $\Omega$
and $253$ heptads. Its point--block incidence matrix is $H$.
Two distinct heptads intersect in one or three points, and $\Gamma$ joins
them in the latter case. Write $A$ for its adjacency matrix. All vector
spaces in this note are over $F=\mathbb F_3$. Put
\[
 E=\{x\in F^{23}:\textstyle\sum_i x_i=0\},\quad
 C=\im A,\quad D=H^{\mathsf T}E.
\]
The shared construction checks the Steiner condition and actual incidence
matrix, not merely the strongly regular parameters.

\begin{theorem}\label{L050.253}
The code $C$ is the orthogonal direct sum $\langle\mathbf j\rangle\oplus D$,
where $D$ is an absolutely irreducible $22$-dimensional module for $M_{23}$.
Both $C$ and $D$ are LCD. Their parameters and minimum multiplicities are
\[
\begin{array}{c|c|r}
\text{code}&\text{parameters}&\text{minimum words}\\\hline
C &[253,23,77]_3&46\\
D &[253,22,97]_3&3542\\
C^\perp &[253,230,6]_3&850080\\
D^\perp &[253,231,6]_3&850080.
\end{array}
\]
The full groups of coordinate permutations preserving $C$ and $\Gamma$
are equal and isomorphic to $M_{23}$.
\end{theorem}

\begin{proof}
The common residual-incidence theorem \cite[Theorem~2.2]{U} applies to the fixed system and gives $HH^{\mathsf T}=2I_{23}$,
$A=H^{\mathsf T}(J_{23}-I_{23})H$, $C=\im H^{\mathsf T}$ and the
orthogonal splitting $C=Fj\perp\!\oplus H^{\mathsf T}(e^\perp)$.
The two point-space summands are nondegenerate because $23\ne0$ in $F$.
Thus it proves both LCD statements and the displayed dimensions.

We next prove the irreducibility, using explicit matrices in place of a
modular character-table inference. The five permutations of $\Omega$
in \texttt{generators.json} preserve the full heptad set. They therefore
act on the code coordinates and on $E$. In the basis
$e_i-e_{23}$, $1\leq i\leq22$, the row matrix for a point permutation
$g$ has entries
\[
 M(g)_{ij}=\mathbf1_{j=g(i)}-\mathbf1_{j=g(23)},\qquad 1\leq i,j\leq22.
\]
An indicator with point $23$ is zero because $j\leq22$.
The certificate checks that multiplication by these matrices agrees
with the induced permutation of all $253$ coordinates on the actual
rows $h_i-h_{23}$. It produces $484$ independent words in these matrices,
each of length at most six, and reevaluates every word before reducing
the flattened matrices. They span $M_{22}(F)$.

The common full-matrix-algebra criterion \cite[Theorem~6.2]{U} proves absolute irreducibility of $E$ and its equivariant image $D$.

The existing enumeration records contain the complete weights of the
row-reduced adjacency code and of the difference code. Their bases are
compared byte for byte with the bases reconstructed here. The latter
row space is $D$: its dimension is $22$, it lies in $C$ and is orthogonal
to $\mathbf j$. The recorded enumerations have respectively $3^{23}$
and $3^{22}$ coefficient vectors. They give the two distances and counts
in the first two rows of the table. The dual distributions follow from the common MacWilliams theorem \cite[Theorem~7.1]{U}.
All coefficients are evaluated over the integers. Nonnegativity,
integrality and total $3^{253-k}$ are checked. Both duals have no words
of weights one through five and exactly $850080$ of weight six.
This proves the table without repeating either exhaustive enumeration.

The $46$ words $h_i,-h_i$ are precisely the minimum words, by their weights and the complete count. The common Steiner recognition proposition \cite[Proposition~2.3]{U} therefore identifies $\Aut(C)$ with the block action of $\Aut(\mathcal D)$. This explicitly discharges its minimum-word hypothesis.

For an explicit upper bound, form the bipartite point--block incidence
graph, with its two parts initially given different colours. Refine the
colours by the old colour and the counts of neighbours of each colour.
After individualizing points $1,2,3,4,5,6$ successively, the certificate
obtains current cell sizes
\[
 23,\quad22,\quad21,\quad20,\quad16,\quad3.
\]
All points are then distinguished. A pointwise automorphism fixes every
heptad, since distinct heptads have different point neighbourhoods.
The common refinement theorem \cite[Theorem~6.4]{U} consequently bounds $|\Aut(\mathcal D)|$ by $23\cdot22\cdot21\cdot20\cdot16\cdot3=10200960$.
Exact group computations on the literal generators give a subgroup of
this order and an isomorphism to the standard group $M_{23}$.
The induced block action is checked to be faithful. Thus the upper bound
is attained and identifies the full design and code groups. Finally the
subgroup preserves $\Gamma$, and every graph automorphism preserves its
adjacency row space. Therefore
$G\leq\Aut(\Gamma)\leq\Aut(C)=G$, proving the asserted equality.
The upper bound uses only incidence refinement; graph-isomorphism
software was used to discover the generators and is not a final premise.
\end{proof}

\paragraph{Qualifications to Proposition 5.}
The main Theorem~1(d) is recovered as stated. Proposition~5 additionally
assigns distance eight to $D^\perp$ and declares it irreducible. Both
assertions are false for these codes. Since $D\subset C$, the invariant
space $C^\perp$ of dimension $230$ is a proper nonzero subspace of
$D^\perp$, whose dimension is $231$. The same inclusion already gives
$d(D^\perp)\leq d(C^\perp)=6$; the exact value six is verified above.
The proposition also calls $D$ the ``unique and smallest'' irreducible
module without qualifications. The unqualified statement overlooks the
trivial one-dimensional module. Its intended classification requires an
input separate from our matrix certificate. We retain the complete
$3$-modular character table of $M_{23}$, as in the source's reference [13],
Jansen--Lux--Parker--Wilson, \emph{An Atlas of Brauer Characters} (1995).
Its irreducible degrees over an algebraic closure are
\[
 1,22,45,45,104,104,231,253,770,770,990,990,1035.
\]
The accompanying GAP calculation checks these entries against its
installed table \texttt{M23mod3}; it does not reprove completeness of
that table. Together with standard scalar extension and descent for
simple modules over finite fields, the table gives a unique
$22$-dimensional irreducible $F$-module, and $22$ is the least
nontrivial degree. Indeed a simple module splits after scalar extension
into an orbit of absolutely simple constituents of equal degree. Below
degree $22$ the only available constituent is trivial, which forces the
original simple module to be trivial. At degree $22$ the unique
constituent is the one already realized by $D$; descent of module
isomorphisms gives uniqueness over $F$. These are retained
classification premises, explicitly distinct from the direct proof of
irreducibility of $D$. The source's open uniqueness question for the
strongly regular parameters $(253,140,87,65)$ is not resolved here.

